\documentclass[11pt]{article}
\usepackage[margin=1in]{geometry}
\usepackage{amsmath,amssymb,amsthm,mathtools}
\usepackage{graphicx}
\usepackage{microtype}
\usepackage{hyperref}
\usepackage{enumitem}
\hypersetup{colorlinks=true,linkcolor=blue,citecolor=blue,urlcolor=blue}

\newtheorem{theorem}{Theorem}
\newtheorem{lemma}{Lemma}
\newcommand{\T}{\mathbb T}
\newcommand{\F}{\mathbb F}
\newcommand{\Rok}{\operatorname{dim}_{\mathrm{Rok}}}
\newcommand{\eind}{\operatorname{ind}}
\newcommand{\genus}{\operatorname{gen}}

\title{An Infinite-Product Counterexample for Rokhlin Dimension\\
\large Candidate negative answers to Questions 5.2 and 5.3 of arXiv:1407.1277}
\author{Automated mathematical discovery run \texttt{fa\_banach\_001}, lane 15}
\date{11 August 2026}

\begin{document}
\maketitle

\noindent\fbox{\parbox{0.95\textwidth}{\textbf{Status.} Candidate full counterexamples, likely valid, requiring expert review. The construction is a constant-order version of the infinite product suggested in the source.}}

\begin{abstract}
Let \(G=\prod_{j\geq1}\mathbb Z/2\mathbb Z\) act on the infinite torus \(X=\prod_{j\geq1}S^1\) by coordinatewise signs. The action is free and the induced action on \(C(X)\) has the \(X\)-Rokhlin property. We prove that its Rokhlin dimension is nevertheless infinite. For the quotient \(H_i=(\mathbb Z/2)^i\), the associated principal bundle squares the first \(i\) circle coordinates. Its mod-2 kernel cup length is at least \(i\), hence its Schwarz genus and equivariant index are at least \(i\). A published inverse-limit formula then makes the Rokhlin dimension the supremum of these indices. Since \(C(X)\) is commutative, the commuting-tower dimension is infinite as well. This gives simultaneous negative answers to Questions 5.2 and 5.3 as stated.
\end{abstract}

\section{The source questions}

Gardella asks whether finite Rokhlin dimension with commuting towers is equivalent to the \(X\)-Rokhlin property for arbitrary compact groups, and whether, in the commutative case, finite Rokhlin dimension is equivalent to freeness of the induced action on the maximal ideal space \cite[Questions 5.2--5.3]{Gardella}. The source immediately suggests an infinite product of cyclic rotations but leaves its Rokhlin dimension unchecked.

\begin{center}
\includegraphics[width=\textwidth]{figures/open_problem_crop.png}

\smallskip
\small Figure 1. Questions 5.2 and 5.3 and the proposed infinite-product route in the source, PDF p.~28.
\end{center}

\noindent\textbf{Transcription.} Question 5.2 asks whether finite Rokhlin dimension with commuting towers is equivalent to the \(X\)-Rokhlin property for arbitrary compact groups. Question 5.3 asks whether finite Rokhlin dimension for actions on commutative \(C^*\)-algebras is equivalent to freeness on the maximal ideal space for arbitrary compact groups.

\section{Proof intuition}

Finite Rokhlin dimension can be tested on the finite Lie quotients of a compact metrizable group. The product group \(G=(\mathbb Z/2)^{\mathbb N}\) has quotients \(H_i=(\mathbb Z/2)^i\). For the coordinatewise action on the infinite torus, the \(H_i\)-bundle seen at level \(i\) is the product of \(i\) nontrivial double covers. Each double cover kills one degree-one mod-2 class, but the product of all \(i\) classes survives on the base. Sectional-category cup length therefore forces at least \(i+1\) local sections, in the convention that the genus is at least \(i\). These lower bounds diverge with \(i\), so the inverse-limit formula forces infinite Rokhlin dimension. Freeness and the \(X\)-Rokhlin property, by contrast, are immediate from the coordinate action and the identity map on \(C(X)\).

\section{The counterexample}

Write \(\mathbb Z_2=\{1,-1\}\). Set
\[
G=\prod_{j=1}^{\infty}\mathbb Z_2,
\qquad
X=\prod_{j=1}^{\infty}\T,
\]
and define
\[
(g_j)_{j\geq1}\cdot(z_j)_{j\geq1}=(g_jz_j)_{j\geq1}.
\]
Both \(G\) and \(X\) are compact metrizable, and the action is continuous.

\begin{theorem}\label{thm:main}
Let \(\alpha:G\to\operatorname{Aut}(C(X))\) be the induced action. Then:
\begin{enumerate}[label=\textup{(\roman*)}]
\item the action of \(G\) on \(X\) is free;
\item \(\alpha\) has the \(X\)-Rokhlin property;
\item \(\Rok(\alpha)=\Rok^{\mathrm c}(\alpha)=\infty\).
\end{enumerate}
Consequently Questions 5.2 and 5.3 of \cite{Gardella} both have negative answers.
\end{theorem}

We first record the elementary cohomological estimate used below. We use the convention of \cite{GHTW}: \(\genus(p)\leq m\) means that the base is covered by \(m+1\) open sets on each of which the bundle admits a section.

\begin{lemma}[Kernel cup-length bound]\label{lem:cup}
Let \(p:E\to B\) be a principal bundle. If there are positive-degree classes
\(v_1,\ldots,v_k\in\ker(p^*:H^*(B;R)\to H^*(E;R))\) with
\(v_1\smile\cdots\smile v_k\neq0\), then \(\genus(p)\geq k\).
\end{lemma}

\begin{proof}
Suppose \(\genus(p)\leq k-1\), and choose a sectional cover \(B=U_1\cup\cdots\cup U_k\). If \(s_r:U_r\to E\) is a section, then
\[
v_r|_{U_r}=s_r^*p^*v_r=0.
\]
Thus \(v_r\) lifts to \(H^*(B,U_r;R)\). The relative cup product of these lifts lies in
\(H^*(B,U_1\cup\cdots\cup U_k;R)=H^*(B,B;R)=0\), so
\(v_1\smile\cdots\smile v_k=0\), a contradiction.
\end{proof}

\begin{proof}[Proof of Theorem \ref{thm:main}]
If \(g\cdot x=x\), then \(g_jx_j=x_j\) for every \(j\). Since \(x_j\in\T\), it follows that \(g_j=1\) for every \(j\), so the action is free.

The \(X\)-Rokhlin property asks for an approximate equivariant central unital homomorphism from \(C(X)\), equivalently an equivariant unital homomorphism into the central sequence algebra \cite[Definition 1.5]{HP}. Here the identity homomorphism \(C(X)\to C(X)\), followed by the constant-sequence inclusion, already has all these properties. Hence \(\alpha\) is \(X\)-Rokhlin.

For \(i\geq1\), let
\[
N_i=\{g\in G:g_1=\cdots=g_i=1\}
   \cong\prod_{j>i}\mathbb Z_2,
\qquad
H_i=G/N_i\cong(\mathbb Z_2)^i.
\]
Then \(N_1\supset N_2\supset\cdots\), \(\bigcap_iN_i=\{e\}\), and \(G=\varprojlim H_i\). Put \(Y_i=X/N_i\) and \(B=X/G\). Using \(\T/\{\pm1\}\cong\T\), with quotient coordinate \(z\mapsto z^2\), identify both \(Y_i\) and \(B\) with \(\T^{\mathbb N}\). Under these identifications the quotient of the free \(H_i\)-action is
\[
q_i:Y_i\longrightarrow B,
\qquad
q_i(z_1,z_2,\ldots)
=(z_1^2,\ldots,z_i^2,z_{i+1},z_{i+2},\ldots).
\]

Let \(\eta\) be the generator of \(H^1(\T;\F_2)\), and let
\(u_j=\pi_j^*\eta\in H^1(B;\F_2)\), where \(\pi_j\) is the \(j\)-th coordinate projection. For \(1\leq j\leq i\), naturality and the degree-two circle map give
\[
q_i^*u_j=2\,\pi_j^*\eta=0
\quad\text{in }H^1(Y_i;\F_2).
\]
On the other hand,
\[
u_1\smile\cdots\smile u_i\neq0.
\]
Indeed, restrict along the inclusion \(\T^i\hookrightarrow\T^{\mathbb N}\) obtained by setting all later coordinates equal to \(1\); the restriction is the nonzero top K\"unneth class of \(\T^i\). Lemma \ref{lem:cup} therefore gives
\[
\genus(q_i)\geq i.
\]

For compact Hausdorff group actions on compact Hausdorff spaces, Schwarz genus equals the equivariant index \cite[Proposition 4.4]{GHTW}. Moreover, the inverse-limit formula of \cite[Corollary 5.15]{GHTW} gives
\[
\Rok(\alpha)
=\sup_i\eind_{H_i}(X\mathbin{\times_G}H_i)
=\sup_i\eind_{H_i}(X/N_i)
=\sup_i\genus(q_i)
\geq\sup_i i
=\infty.
\]
Thus \(\Rok(\alpha)=\infty\). Finally, all ranges in the central sequence algebra of the commutative algebra \(C(X)\) commute, so ordinary and commuting-tower Rokhlin dimensions agree. Hence \(\Rok^{\mathrm c}(\alpha)=\infty\) as well.

The free action with infinite Rokhlin dimension disproves the equivalence in Question 5.3. The same action is \(X\)-Rokhlin but has infinite commuting-tower Rokhlin dimension, disproving the equivalence in Question 5.2.
\end{proof}

\section{Verification, scope, and novelty}

The proof audit checked compact metrizability, continuity and freeness, the exact \(X\)-Rokhlin definition, the finite-quotient bundle coordinates, the relative-cohomology cup-length argument, equality of genus and index, and every hypothesis of the inverse-limit formula. No computation is used. The supporting formula occurs as Theorem 5.14 and Corollary 5.15 of \cite{GHTW}; that paper does not appear to compute this product example.

The source suggested \(G=\prod_n\mathbb Z_n\) acting on \(\prod_nS^1\). We use the cleaner constant-order group \(\prod_j\mathbb Z_2\), for which arbitrarily long mod-2 cup products give a uniform lower-bound mechanism. This settles the two yes/no questions but does not address the source's other open problems.

A bounded search on 11 August 2026 covered all four run indexes, exact question wording, the arXiv id and title, the published source, the author's thesis, \(X\)-Rokhlin and arbitrary-compact-group phrases, coordinatewise rotations, infinite Rokhlin dimension, Schwarz genus, arXiv:1801.00767, and later citations returned by those searches. It found the original suggestion and the later general formula but no explicit calculation of this example or announced resolution of Questions 5.2--5.3. Novelty is therefore provisional.

\paragraph{Human-review recommendation.}
Check the identification of \(q_i\), the genus convention in Lemma \ref{lem:cup}, and the application of Corollary 5.15 of \cite{GHTW}. Also confirm that Question 5.2 uses the \(X\)-Rokhlin property with an arbitrary compact free \(G\)-space, as in \cite{HP}; the underlying \(X\) then supplies the required witness.

\begin{thebibliography}{9}
\bibitem{Gardella}
E.~Gardella,
\emph{Rokhlin dimension for compact group actions},
Indiana Univ. Math. J. \textbf{66} (2017), no.~2, 659--703;
arXiv:1407.1277.

\bibitem{GHTW}
E.~Gardella, P.~M.~Hajac, M.~Tobolski, and J.~Wu,
\emph{The local-triviality dimension of actions of compact quantum groups},
arXiv:1801.00767v3 (2019). See Proposition~4.4, Theorem~5.14, and Corollary~5.15.

\bibitem{HP}
I.~Hirshberg and N.~C.~Phillips,
\emph{Rokhlin dimension: obstructions and permanence properties},
Doc. Math. \textbf{20} (2015), 199--236, especially Definition~1.5;
arXiv:1410.6581.
\end{thebibliography}

\end{document}
